\chapter{What Matters: A Global Analysis}

The preceding chapters each tested a specific axis group in isolation.
This chapter inverts that framing and treats the entire experiment
database, all models trained across Chapters 4--6 plus the existing
exploratory runs, as the object of study. The question is global:
across everything trained, which architectural axes consistently move
the needle, and which have no measurable effect?

No new models are trained here. At the end of all training the W\&B
database is frozen, all runs are backfilled with a standardised axis
vector using the V2 reference, and the result is a single clean table:
one row per model, one column per axis, plus AUROC and efficiency
metrics. All analysis in this chapter operates on that table.

% ──────────────────────────────────────────────
\section{The run database}

After backfilling, every run carries a complete axis vector covering
T1, T1-a, T1-b, D01--D03, E1, A1--A6, F1, B1 and its active
sub-axes, C1, and the model-size keys H01--H04. Runs are filtered to
the primary task (4t vs.\ combined background) and the default model
size to keep AUROC values on a comparable scale. The resulting table
contains roughly $N$ models and forms the sole data source for this
chapter.

% ──────────────────────────────────────────────
\section{Marginal sensitivity analysis}

For each axis, the conditional distribution $P(\text{AUROC} \mid
\text{axis} = s_i)$ is estimated by grouping all runs sharing setting
$s_i$ and computing the group mean, marginalising over all other axes.
Because sweeps are approximately orthogonal, these marginal means are
interpretable as main effects. The primary summary statistic per axis
is the range: the difference between the highest and lowest group mean.
Axes where this range falls below the seed-level noise floor
($\lesssim 0.002$ AUROC) are classified as axes that do not matter.

The main output is a single ranked bar chart spanning all axes,
colour-coded by family (D, T, E, A, F, B, C). This is the headline
figure of the chapter. Selected null results, including PE type, norm
policy, and causal masking, are each given a small panel showing the
flat marginal distribution, making the null result explicit rather than
just asserted.

% ──────────────────────────────────────────────
\section{Surrogate modelling and configuration prediction}

Beyond marginal effects, a gradient-boosted decision tree is trained
on the run database to predict AUROC from the full axis configuration
vector. This captures non-linear interactions between axes that the
marginal analysis cannot detect, for example a bias family that only
helps when combined with the identity tokeniser.

The idea is simple: treat every trained model as a data point, with
its axis settings as features and its AUROC as the label. Once the
tree is fitted, it can be queried on axis combinations that were never
trained, giving a predicted AUROC for any configuration in the full
space. Feature importances from the tree provide a second, independent
ranking of axis sensitivity that can be compared directly to the
marginal analysis above.

The top predicted configurations are reported as hypotheses, not
established results, and feed directly into the recommended model of
Ch.\,9. No new training is required.

% ──────────────────────────────────────────────
\section{Recommended configuration}

The findings of the marginal analysis and surrogate model are combined
into a single recommended axis configuration: each axis is set to
whichever value shows a robust improvement, or to the cheapest default
where no improvement was found. The recommended configuration is stated
explicitly as a Hydra override string and validated in Ch.\,9.
